\section{Complete Phase 5B Implementation Details}
\label{supp:implementation}

\subsection{Conditions and budgets}

Phase~5B is a matched ablation over the same 18 cells.
\begin{table*}[t]
\centering
\begin{tabular}{@{}clllrr@{}}
\toprule
 & Implementation & Workflow & Instructions & Calls & Max/call \\
\midrule
A & M-short-2 & monolithic & concise & 2 & 4096 \\
B & M-unified-2 & unified & enriched & 2 & 4096 \\
C & M-unified-4 & unified staged & same as B & 4 & 2048 \\
D & Specialized-4 & four roles & role-specific & 4 & 2048 \\
\bottomrule
\end{tabular}
\caption{Frozen Phase~5B conditions and nominal output budgets.}
\end{table*}
Every condition has nominal output capacity 8,192 tokens.  B and C use the
same substantive four-role instruction content.  B uses one persona and two
calls; C uses one continuing workflow and four sequential calls.  D uses
specialized planner, provenance, reviewer, and reporter prompts and differs
from C in message/context partitioning and intermediate interfaces.

The preregistered primary contrast \(C-B\) changes the number and organization
of calls while holding substantive instructions and nominal output capacity
fixed.  Secondary contrast \(B-A\) tests instruction enrichment under
two-call implementations.  \(D-C\) is a residual architectural comparison
bundling role scope, context, interfaces, and message flow.  The study is a
controlled component ablation, not a complete factorial causal decomposition.

\subsection{Accounting and analysis}

The design contains 72 intended units and 216 planned calls:
\(18(2+2+4+4)\).  The scheduler rotates condition order by cell.  Provider
retries and schema repairs are zero.  Planned calls are not treated as
observed dispatches.

Point estimates are paired all-cell mean regret differences.  Supporting
analyses use completed pairs, scenario-cluster bootstrap intervals,
scenario-family and seed subsets, the prespecified diagnostic subset, and
leave-one-extreme-cell trimming.  With only \(K=2\) scenario clusters, all
claims remain pilot-descriptive; intervals are not population-level
significance tests \citep{MacKinnonWebb2018}.
